\section{Introducción: por qué los datos vuelven a ser el eje de la industria}

Suele decirse que la inteligencia artificial se impulsa con tres motores: los datos, la capacidad de cómputo y los algoritmos. En los últimos años, la capacidad de cómputo y la arquitectura de los modelos han sido más visibles; pero ahora que los modelos de lenguaje de gran tamaño se acercan al «muro de datos» y la robótica enfrenta un «desierto de datos», los datos vuelven a ser la variable que determina el panorama de la industria. El valor de esta entrevista con Xie Chen (谢晨) es que no habla de un solo conjunto de datos, sino que descompone la industria de los datos en fuentes, ciclo cerrado, pirámide, fijación de precios, Recipe y mapa del sector.

\begin{importantbox}{Tesis central de este episodio}
Los datos no son una materia prima estática, sino un ciclo cerrado continuo. Quien obtenga a menor costo fallas reales, escenarios simulados, retroalimentación verificable y recetas reutilizables podrá entrenar más rápido sistemas inteligentes que se puedan llevar a la práctica.
\end{importantbox}

\begin{knowledgebox}{Nota sobre la estrategia visual}
Este video es una entrevista con encuadre fijo, sin slides didácticas, pizarrón ni demostraciones de producto. Conforme al estándar de este repositorio para pódcasts, el cuerpo del texto no repite fotogramas de los participantes; la portada sirve para identificar la fuente, y el contenido didáctico se presenta mediante la pirámide de datos, el ciclo cerrado y el mapa de la industria.
\end{knowledgebox}

\subsection{Resumen de la sección}

Este episodio es un panorama de la industria de los datos: el problema de los LLM es el rendimiento marginal decreciente de los datos de alta calidad; el de la robótica, que los datos del mundo físico son escasos y costosos. Ambos necesitan un ciclo cerrado de datos y un sistema de evaluación.
